\begin{figure}[!t]
    \centering
    \includegraphics[width=1.0\linewidth]{figs/caption_demo5.png} 
    \vspace{-6mm}
    \caption{The illustration of output textual caption types and intermediate texts for caption selection.
    %The output encompasses short, long, and structured captions specifically designed to address human characters. 
    %The intermediate results include the original captions generated by CogVLM and MiniCPM, and the recaptioned texts employed in the selection process. 
    The deep red text represents key information.}
    \vspace{-7mm}
    \label{fig:caption_demo} 
\end{figure}

\begin{figure*}[!t]
    \centering
    \begin{subfigure}{0.33\linewidth}
        \centering
        \includegraphics[width=0.95\textwidth]{figs/filtered_video/subject_consistency_clipped.png} 
        \vspace{-1mm}
        \caption{Subject consistency ($\uparrow$)}  
        % \vspace{1mm}
        \label{fig:Vbench metirc}  
    \end{subfigure}%
    \hfill
    \begin{subfigure}{0.33\linewidth}
        \centering
        \includegraphics[width=0.95\textwidth]{figs/filtered_video/dynamic_degree_clipped.png}  
        \vspace{-1mm}
        \caption{Dynamic degree}  
        % \vspace{1mm}
        \label{fig:Vbench metirc}  
    \end{subfigure}%
    \hfill
    \begin{subfigure}{0.33\linewidth}
        \centering
        \includegraphics[width=0.95\textwidth]{figs/filtered_video/aesthetic_quality.png}  
        \vspace{-1mm}
        \caption{Aesthetic quality ($\uparrow$)} 
        % \vspace{1mm}
        \label{fig:Vbench metirc}  
    \end{subfigure}  
    \begin{subfigure}{0.33\linewidth}
        \centering
        \includegraphics[width=0.95\textwidth]{figs/filtered_video/imaging_quality.png}  
        \vspace{-1mm}
        \caption{Image quality ($\uparrow$)}
        % \vspace{1mm}
        \label{fig:Vbench metirc}  
    \end{subfigure}%
    \hfill
    \begin{subfigure}{0.33\linewidth}
        \centering
        \includegraphics[width=0.95\textwidth]{figs/filtered_video/temporal_flickering.png}  
        \vspace{-1mm}
        \caption{Temporal flickering ($\uparrow$)}  
        % \vspace{1mm}
        \label{fig:Vbench metirc}  
    \end{subfigure}%
    \hfill
    \begin{subfigure}{0.33\linewidth}
        \centering
        \includegraphics[width=0.95\textwidth]{figs/filtered_video/motion_smoothness.png}  
        \vspace{-1mm}
        \caption{Motion smoothness ($\uparrow$)}  
        % \vspace{1mm}
        \label{fig:Vbench metirc}  
    \end{subfigure}  
    \vspace{-4mm}
    \caption{The comparison of video quality between Panda-70M and the proposed data before and after video quality filtering.
     We utilize the video quality evaluation metrics introduced in VBench to assess the video quality. 
     We can see that the general video quality of the proposed data has obviously improved after video quality filtering and superior to that of the Panda-70M.}
     \label{fig:quality}
     \vspace{-5mm}
\end{figure*}

\section{Dataset}
% This section begins with a discussion of the source dataset in Section~\ref{sec:source_data}. 
% Section~\ref{sec:data_process} outlines the data processing steps, the generation of human motion, and a quality filter that considers the alignment between textual captions and human appearance, human motion, and face motion, as well as the synchronization between speech audio and human talking video. 

\subsection{Source Data}\label{sec:source_data}
As shown in Figure~\ref{fig:source_data}, the dataset comprises a total of 134,000 audiovisual works, with an aggregate runtime of 105,000 hours. This collection includes 5,192 films, 4,289 television series, and 192 documentaries, spanning a diverse range of genres across 15 categories, including Action, Adventure, Animation, Comedy, Crime, Documentary, Drama, Fantasy, Kids, Mystery, Politics, Romance, Science Fiction, Soap and War. The temporal scope of the dataset extends from the 1920s to the present day, featuring content in 58 different languages.
Table~\ref{tab:datasets} presents a comparative analysis of our dataset against previous general and human video datasets. The proposed human-centric dataset achieves a scale comparable to several well-known general video datasets while incorporating motion conditions typically found in human-specific datasets.
In terms of technical specifications, 84.5\% of the videos possess a resolution exceeding 720P, while 52.7\% exceed 1080P. 
% Furthermore, 74.0\% of the data exhibits a frame rate greater than 24 FPS, with an average bitrate surpassing 3,500 kbps. 

We selected films, television series, and documentaries as our primary data sources due to their emphasis on character-driven narratives, which typically exhibit superior luminance, aesthetic quality, and motion characteristics. 
Figure~\ref{fig:quality} demonstrates that the video quality after preprocessing is superior when compared to the Panda-70M dataset, which is primarily sourced from the internet.

\subsection{Data Processing}\label{sec:data_process}
\begin{figure}[!t]
    \centering
    \includegraphics[width=1.0\linewidth]{figs/Pipeline.png} 
    \vspace{-6mm}
    \caption{Overview of the proposed extended DiT-based video generation models.}
    \label{fig:network_pipeline}
    \vspace{-7mm}
\end{figure}

\subsubsection{Video Preprocessing}
The initial phase of video preprocessing encompasses several essential operations aimed at enhancing data quality. 
This phase includes codec standardization using H.264 and subtitle cropping via the CRAFT method~\cite{baek2019character}, which effectively removes spatial regions associated with subtitles.
Additionally, black border cropping is performed by analyzing pixel values across each frame and row to eliminate unwanted borders. Scene splitting is executed using SceneDetect~\cite{PySceneDetect}, which assesses color and brightness variations to identify significant transitions and determine potential split points. Finally, a duration filter is applied to trim video segments to a specified length of 2 to 20 seconds, ensuring the retention of relevant content for further analysis.

\subsubsection{Video Quality Filter}

In this study, we implement a multi-faceted video quality filtering process based on several key metrics: luminance, blurriness, aesthetic appeal, global and local motion, and overall technical quality. Each metric is evaluated against specific thresholds to ensure that only high-quality video content is retained for further processing.
Please see Figure~\ref{fig:quality} for the comparison of video quality before and after this filtering process.

\subsubsection{Human Motion Generation}
\label{sec:human_motion}
\noindent\textbf{Textual Caption.}
% In video generation models, textual captions play a critical role by providing contextual information that guides the generation process, enhances multimodal learning, and improves model generalization. 
% They encourage alignment between generated visual content and textual descriptions.
As illustrated in Figure~\ref{fig:caption_demo}, we employ two publicly available multimodal models, namely MiniCPM~\cite{yao2024minicpm} and CogVLM~\cite{hong2024cogvlm2}, to generate textual captions for the provided videos. 
MiniCPM has been extensively utilized in prior research~\cite{chen2024multi}, while CogVLM is a more recent introduction to the field. 
To enhance the quality of the generated captions, we implement a voting strategy. 
However, the limited instruction-following capabilities of these models often lead to captions with chaotic structures that lack a unified format, thereby complicating the training of video generation models. 
To address this issue, we utilize LlaMA 3.1~\cite{dubey2024LlaMA} to reformat the initial captions. 
Subsequently, we employ BLIP2~\cite{li2023blip} to select captions with higher similarity scores from the two generated sets.

The final captions are provided in short, long and structured formats to ensure a comprehensive representation of the video content. 
Long captions are derived from the aforementioned process. 
To obtain structured captions, we utilize LlaMA to extract components related to human appearance, human motion, face motion and environment description. 
Finally, we obtain the short captions by merging the structured captions by LlaMA and restricting caption length to approximately 20 to 30 words.





\noindent\textbf{Skeleton Sequence.}  
The adoption of skeleton poses has become prevalent in diffusion-based human video generation models~\cite{zhu2024champ,chang2023magicpose,hu2024animate}. 
In our dataset, we utilize DWpose to extract skeleton data, which serves as the foundational annotation for controlling human poses in the generated videos. 
The human skeleton sequences obtained through this method, along with the corresponding human videos, aim to enhance the existing human animation datasets by contributing to the diversity of characters, scenes, and movements.


\noindent\textbf{Speech Audio.}  
Recent studies have investigated audio-driven head animation ~\cite{xu2024hallo,xu2024vasa,he2023gaia}, gesture generation ~\cite{yang2023qpgesture,zhu2023taming}, and holistic human motion ~\cite{corona2024vlogger,wang2024dance}, underscoring the significance of audio as a guiding modality. 
In our dataset, we include the corresponding audio tracks for each video clip. 
Furthermore, we employ SyncNet~\cite{raina2022syncnet} to assess the consistency between the speech audio of the subject and their lip motion. 
We anticipate that the integration of speech audio in human videos will further advance research in audio-driven animation, extending its applicability to more realistic and generalized scenarios.





\subsubsection{Human Quality Filter}
\label{sec:caption_refine}

This phase focuses on assessing the presence and characteristics of human subjects within the video. Key criteria for filtering include the number of individuals detected, the dimensions of human and facial regions, and the velocity of human motion. Additionally, the alignment of text and video concerning human appearance, motion, and facial expressions is also considered.






\noindent\textbf{Human Appearance Text-Video Alignment.}
We assess the alignment between the generated video and the human appearance descriptions provided in the text prompt, which include characteristics such as physical features, style, clothing, and accessories. 
By extracting the appearance-related text from the generated long caption, we compute the BLIP2 score~\cite{li2023blip} to evaluate the correspondence between the appearance text and the generated video. 
This process ensures that the visual attributes of the human subjects are accurately represented in accordance with the textual descriptions.






\noindent\textbf{Human Motion Text-Video Alignment.}
We evaluate the alignment of the generated video with the human motion described in the text prompt, which encompasses aspects such as action type, speed, and trajectory. 
Utilizing the LlaMA model, we extract the action-related text from the original caption and compute the BLIP2 score~\cite{li2023blip} to assess the semantic alignment between the action text and the generated video. 
This evaluation promotes the accurate depiction of actions in the video as intended by the descriptions.

\noindent\textbf{Face Motion Text-Video Alignment.}
Similarly, we analyze the extent to which the generated video aligns with facial motion, including expressions, head posture, and emotional variations described in the text prompt. 
By extracting the emotion-related text using the LlaMA model, we calculate the BLIP2 score~\cite{li2023blip} to ensure that the emotional content is accurately conveyed. 
This alignment is essential for maintaining the intended emotional narrative of the video.